\section{Limitations}
\label{sec:limitations}

% Moved from the report. State each weakness and why it does not undermine the main claims.

\textbf{Three datasets, chosen for contrast.} The main study uses CIFAR-100, whose $32\times32$ images are upsampled to $224\times224$ for the classifiers. Stage 3 adds Flowers-102 and HAM10000 at $224\times224$, and there the main CIFAR-100 result does not hold. The two datasets differ from CIFAR-100 in more than their distance from Stable Diffusion: in the number of classes (102 and 7), class balance, image content and metric. With one dataset per kind, we cannot say which of these differences causes the reversal; the sample grids point to class fidelity and the domain gap, but that is an observation, not a measured cause.

\textbf{Seeds.} The comparisons behind the main claims have five training seeds, the others three (Section~\ref{sec:method}). A gain or loss is declared only when every seed agrees in sign, which with three seeds happens by chance a quarter of the time when there is no effect, and the standard deviations of so few seeds are themselves uncertain. The paired spread is largest from random initialization: on CIFAR-100 it reaches 5.7 points (ResNet-50, best generated-image condition against best conventional at 50 images per class), and on HAM10000 standard deviations of 2 to 5 points are common. The study can therefore resolve only effects larger than this spread; smaller ones are reported as inconclusive, and HAM10000 has many.

\textbf{Fixed step budgets.} The step budgets were chosen in pilots on the real-only condition and are equal across conditions, so the comparison holds the training compute fixed. Some runs reached their best validation score at the final step: in the pretrained arm 1 of 30 real-only and 8 of 30 RandAugment runs, and in the random-initialization arm 6 of 30 and 13 of 30. Because the cosine schedule ends at a learning rate of zero, the last steps are often the best, but augmented runs might also gain from longer training, so the RandAugment gains may be understated. In the synthetic-data conditions the same step budget is spread over twice as many images, so each real image is seen about half as often as in the real-only run; the exposures are recorded for every run.

\textbf{One fixed recipe from scratch.} Absolute accuracy from random initialization is low: the best CIFAR-100 result on the full training set is 65\% (ResNet-50 with RandAugment), far below what long, tuned schedules reach. Every condition of a budget shares one recipe and one step budget, chosen in pilots on real-only training, because the question is what generated images add at a fixed training cost; tuning each condition separately would mix the effect of the images with the effect of the tuning. The cost is that the comparison holds for this recipe. With longer training, real-only and RandAugment runs may gain more than runs that already see more varied images, or less, and we have not measured which.

\textbf{Outside pretraining.} The ImageNet weights of the pretrained classifiers and the web-scale training data of Stable Diffusion are outside our control. The diffusion training data may contain images similar to, or taken from, CIFAR-100, so any benefit of generated images may partly reflect that exposure. The random-initialization arm removes the classifier's outside knowledge, and the class-conditional GAN, trained only on the 5{,}000 real images of one budget, removes the generator's; but the GAN was run at that one budget only.

\textbf{Two software environments.} The synthetic-data runs with seeds 0 to 2 were trained with newer Python, PyTorch, Torchvision and NumPy versions than the real-only and RandAugment runs they are paired with (Section~\ref{sec:setup}). For seeds 3 and 4, both sides of each pair used the same versions. The recipe, split and step budgets are identical, but a version change can shift accuracy slightly, so small paired changes between the two groups should be read with that in mind.
% TODO: if the baseline rerun in the newer environment is done, give the measured shift.

\textbf{Stage 3 step budgets and learning rates.} Stage 3 keeps the CIFAR-100 step budget of each images-per-class budget. With 7 classes, HAM10000 has about one-fourteenth as many training images as CIFAR-100 at the same budget, so each image is seen about 14 times as often (a pretrained run at 5 images per class passes over its 35 images about 2{,}200 times). Repetition on this scale may favour the augmented conditions over real-only training; all conditions of a cell share the same budget, so the comparisons between augmentations are not affected in this way. Two pretrained HAM10000 learning rates ($10^{-3}$ for ResNet-50 and $10^{-4}$ for ViT-B/16) and the pretrained ViT-B/16 rate for Flowers-102 ($2\times10^{-4}$) were the highest rates tried, so a higher rate might have done better for every condition of those arms. The HAM10000 pilots were not extended: at 20 images per class (seed 100, validation balanced accuracy), ResNet-50 reached 56.3, 58.6 and 61.3\% at $10^{-4}$, $3\times10^{-4}$ and $10^{-3}$, and ViT-B/16 reached 56.6, 57.9 and 58.8\% at $2\times10^{-5}$, $5\times10^{-5}$ and $10^{-4}$, so both were still rising at the top of the range. The rates from random initialization (1.0 and $3\times10^{-4}$) were inside the range tried.

\textbf{A follow-up chosen after the results.} The Stage 3 LoRA experiment was decided after the Stage 3 test results were seen, to test an explanation of them. Its settings were copied from CIFAR-100 and its memorization rule fixed before any pool was generated, but it is reported as a follow-up, not as part of the planned protocol.
